\section{The dimension threshold for finite semidefinite lifts}
\label{repr:section}

The three-variable hull has a finite semidefinite lift, whether or not the
inequalities studied above give a complete description. We now prove that this
changes at four variables. The obstruction concerns every finite lift, with
arbitrary auxiliary variables and matrix sizes.

Recall the definition of a spectrahedral shadow in \eqref{eq:set:shadow}.
Linear images, linear preimages, affine sections, products, and
Minkowski sums preserve this property. A closed convex cone is a spectrahedral
shadow if and only if its dual is one
\citep[Remark~3.17]{BodirskyKummerThom2024}.

\subsection{Normalization and the positive diagonal directions}

The constant coefficient of a quadratic must be retained when taking a dual.
For a compact convex set $K$ and a polyhedral cone $D$, write
\begin{equation}\label{repr:homogenization}
 \widehat{K+D}
 =\{(t,tk+d):t\geq0,\ k\in K,\ d\in D\}.
\end{equation}
This is the closed homogenized cone; its height-zero section is
$\{0\}\times D$.

\begin{lemma}\label{repr:normalization}
Let $K$ be nonempty, compact, and convex, and let $D$ be a polyhedral cone.
Then $K+D$ has a finite semidefinite lift if and only if
$\widehat{K+D}$ has one.
\end{lemma}
\begin{proof}
The cone in \eqref{repr:homogenization} is closed: in a convergent sequence
$(t_j,t_jk_j+d_j)$, compactness bounds $k_j$, and therefore $d_j$ is bounded;
take convergent subsequences of $k_j$ and $d_j$. Its height-one section is
$K+D$, so one implication follows by taking an affine section.

Conversely, homogenize a lift \eqref{eq:set:shadow} of $S=K+D$ to
\[
 t\geq0,\qquad
 tA_0+\sum_i y_iA_i+\sum_j z_jB_j\succeq0.
\]
At $t>0$ the visible point is precisely $(t,y)$ with $y\in tS$.
If $(0,y,z)$ is feasible, add any nonnegative multiple of it to a feasible
lift $(1,s,z_s)$. This shows $s+\alpha y\in S$ for every $\alpha\geq0$.
Consequently $y$ is a recession direction of $S$. The recession cone of
$K+D$ is $D$: writing $s+\alpha y=k_\alpha+d_\alpha$ and dividing by
$\alpha\to\infty$ proves $y\in D$, while the opposite inclusion is immediate.
Thus the homogenized lift has no height-zero direction outside $D$.
Its projection may omit some directions in $D$; adding the polyhedral cone
$\{0\}\times D$ fills them. At positive height this addition leaves $tS$
unchanged because $S+D=S$. The resulting projection is exactly
\eqref{repr:homogenization}.
\end{proof}

Apply the lemma to $\mathcal Q_n$ and $\mathcal H_n=\mathcal Q_n+D_n$, where
$D_n=\{(0,\diag(s)):s\geq0\}$. Under the evaluation pairing for quadratics,
\begin{equation}\label{repr:duals}
 (\widehat{\mathcal Q_n})^*=\mathcal P_n,
 \qquad
 (\widehat{\mathcal H_n})^*=\mathcal P_n^+.
\end{equation}
Indeed, evaluation on $(1,x,xx^\top)$ requires nonnegativity on $C_n$, and
evaluation on the additional diagonal rays requires $b_i\geq0$.
This identifies both dual cones without discarding the constant coordinate
or the recession directions.

\subsection{Finite lifts through dimension three}

For completeness, we give the established construction that supplies the
positive side of the threshold \citep{AnstreicherBurer2010}. Let
$\CP_r=\cone\{vv^\top:v\in\R_+^r\}$ and
$\DNN_r=\{W\in\Sn{r}:W\succeq0,\ W\geq0\text{ entrywise}\}$.
The classical identity $\CP_r=\DNN_r$ holds for $r\leq4$
\citep{Diananda1962,MaxfieldMinc1962}.

Triangulate $C_n$ into its $n!$ order simplices. If $V_\pi$ is the matrix
of the $n+1$ vertices of simplex $\pi$, introduce
$W_\pi\in\DNN_{n+1}$ and impose
\begin{equation}\label{repr:small-lift}
 \sum_\pi \mathbf1^\top W_\pi\mathbf1=1,
 \qquad m=\sum_\pi V_\pi W_\pi\mathbf1,
 \qquad Y=\sum_\pi V_\pi W_\pi V_\pi^\top.
\end{equation}
For $n\leq3$ this is an exact lift of $\mathcal Q_n$.
To verify exactness, decompose
$W_\pi=\sum_\ell v_\ell v_\ell^\top$ with $v_\ell\geq0$.
Each nonzero summand represents the simplex point
$V_\pi v_\ell/(\mathbf1^\top v_\ell)$ with mass
$(\mathbf1^\top v_\ell)^2$. These masses sum to one in
\eqref{repr:small-lift}. Conversely, each simplex atom has such a rank-one
representation. Adding $D_n$ gives a lift of $\mathcal H_n$.
The same argument applies to the quadratic moment hull of every polytope of
dimension at most three, after triangulation in its affine hull.

\subsection{The four-variable obstruction}

The proof uses completely positive maps on real closed fields, following
\citet{BodirskyKummerThom2024}. Appendix~\ref{repr:appendix} states the precise
external inputs and explains why such maps preserve a finite semidefinite
lift. The needed consequence is as follows. Let $H$ be the Horn matrix, with
diagonal entries $1$, entries $-1$ on the cycle
$12,23,34,45,51$, and entries $1$ at the remaining off-diagonal positions.
There are a real closed field $F\supseteq\R$ and a function
$f:\Q^5\to F$ such that $f(0)=1$,
\begin{equation}\label{repr:negative-horn}
 \eta:=\sum_{i,j=1}^5 H_{ij}f(2e_i+2e_j)<0,
\end{equation}
and the coefficient map
\begin{equation}\label{repr:coefficient-map}
 L_f\left(\sum_a c_a\varepsilon^a\right)
 =\sum_a f(a)c_a\varepsilon^a
\end{equation}
is unital and completely positive on the real closed Hahn field
$E=F((\varepsilon^{\Q^5}))$. Order $\Q^5$ lexicographically, with its first
coordinate most significant. A monomial with positive leading exponent is
then smaller than every positive real constant.

\begin{theorem}\label{repr:threshold}
The sets $\mathcal Q_n$ and $\mathcal H_n$, and the cones
$\mathcal P_n$ and $\mathcal P_n^+$, are spectrahedral shadows if and only if
$n\leq3$.
\end{theorem}
\begin{proof}
The preceding construction, Lemma~\ref{repr:normalization}, and closed-cone
duality prove the assertion for $n\leq3$.

For $n=4$, put $u_1=1$ and $u_{k+1}=x_k$, and consider the quadratic over $E$
\begin{equation}\label{repr:infinitesimal-quadratic}
 q(x)=\sum_{i,j=1}^5 H_{ij}
       \varepsilon^{2e_i+2e_j}u_i u_j.
\end{equation}
The Horn matrix is copositive. This remains true over $E$ by transfer for real
closed fields, so $q(x)\geq0$ for $x\in C_4(E)$. Its square coefficients are
$\varepsilon^{4e_2},\ldots,\varepsilon^{4e_5}>0$, and hence
$q\in\mathcal P_4^+(E)$.

Apply $L_f$ to the coefficients of $q$ and evaluate at
\[
 \xi_k=\varepsilon^{2e_1-2e_{k+1}},\qquad k=1,\ldots,4.
\]
Every $\xi_k$ is positive and smaller than $1$ by the chosen lexicographic
order. Moreover
$\varepsilon^{2e_i}u_i(\xi)=\varepsilon^{2e_1}$ for all $i$.
Therefore
\begin{equation}\label{repr:evaluation}
 L_f(q)(\xi)=\varepsilon^{4e_1}
       \sum_{i,j=1}^5H_{ij}f(2e_i+2e_j)
       =\varepsilon^{4e_1}\eta<0.
\end{equation}
If either $\mathcal P_4$ or $\mathcal P_4^+$ had a finite semidefinite lift,
the same lift would describe that cone over $E$, and the unital completely
positive map $L_f$ would preserve it. Equation~\eqref{repr:evaluation}
contradicts this preservation. Neither cone is a shadow.

For $n>4$, restrict to quadratics depending only on $x_1,\ldots,x_4$.
This linear section of $\mathcal P_n$, respectively $\mathcal P_n^+$, is
$\mathcal P_4$, respectively $\mathcal P_4^+$, so the larger cones cannot
be shadows. Finally, \eqref{repr:duals} and
Lemma~\ref{repr:normalization} rule out lifts of $\mathcal Q_n$ and
$\mathcal H_n$ for all $n\geq4$.
\end{proof}

The constant coordinate in \eqref{repr:infinitesimal-quadratic} supplies the
fifth copositive variable. The evaluation point places the resulting
obstruction in an infinitesimal neighborhood of a cube vertex. The failure
of $\CP_5=\DNN_5$ alone would only rule out the construction
\eqref{repr:small-lift}; Theorem~\ref{repr:threshold} rules out every finite
semidefinite construction. This conclusion is consistent with exact separation
methods for low-dimensional box cones \citep{BurerDong2013}: a separation
procedure need not give a finite lift.
The theorem concerns a single exact formulation of the joint moment hull.
Particular objective epigraphs may have simpler lifts, and convergent
semidefinite hierarchies remain possible. At fixed dimension, the existence
of an exact separation algorithm and the absence of a finite lift are
compatible statements; the theorem makes no algorithmic hardness claim.

\subsection{Simple vertices of other polytopes}

The argument depends on a local orthant at one vertex. The following form
states the geometric condition directly. For a closed convex semialgebraic
cone $K\subseteq\R^m$, define
\[
 \COP(K)=\{M\in\Sn{m}:z^\top Mz\geq0\text{ for every }z\in K\},
 \qquad
 \CP(K)=\cl\cone\{zz^\top:z\in K\}.
\]

\begin{theorem}\label{repr:apex}
Suppose that a linear map $\ell=(\ell_1,\ldots,\ell_5):\R^m\to\R^5$
and a real $\delta>0$ satisfy
\begin{equation}\label{repr:apex-condition}
 \ell(K)\subseteq\R_+^5,
 \qquad
 \{u\geq0:u_k\leq\delta u_1\ (k=2,\ldots,5)\}\subseteq\ell(K).
\end{equation}
Then neither $\COP(K)$ nor $\CP(K)$ is a spectrahedral shadow.
\end{theorem}
\begin{proof}
Use the fields and map in \eqref{repr:negative-horn}--\eqref{repr:coefficient-map}.
The form
\[
 M(z)=\sum_{i,j=1}^5H_{ij}\varepsilon^{2e_i+2e_j}
                         \ell_i(z)\ell_j(z)
\]
is nonnegative on $K(E)$, since $\ell(z)\geq0$ there.
The inclusion in \eqref{repr:apex-condition} is a first-order statement with
real parameters and therefore remains valid over $E$. In particular there
is $w\in K(E)$ with
\[
 \ell(w)=(1,\varepsilon^{2e_1-2e_2},\ldots,
                    \varepsilon^{2e_1-2e_5}).
\]
Because $L_f$ is real-linear, it acts termwise on the real coefficients of
$\ell_i\ell_j$. Consequently
$L_f(M)(w)=\varepsilon^{4e_1}\eta<0$.
Preservation under $L_f$ contradicts a finite lift of $\COP(K)$.
The identity $\CP(K)^*=\COP(K)$ and closed-cone duality give the second
assertion.
\end{proof}

\begin{corollary}\label{repr:simple-vertex}
Let $K$ be a full-dimensional pointed polyhedral cone of dimension at least
five with a simple extreme ray, meaning a ray contained in exactly
$\dim K-1$ facets. Then $\COP(K)$ and $\CP(K)$ are not spectrahedral
shadows. If a polytope $X$ has dimension at least four and a simple vertex,
then its cone of nonnegative quadratics and its quadratic moment hull have
no finite semidefinite lifts.
\end{corollary}
\begin{proof}
At a simple extreme ray, independent incident facet forms give nonnegative
coordinates transverse to the ray. A form strictly positive on the cone
gives the first coordinate. Sufficiently small nonnegative transverse
coordinates satisfy the other, strictly slack facet inequalities. This
proves \eqref{repr:apex-condition}; Appendix~\ref{repr:simple-ray-proof}
gives the uniform bound and covers height zero.

For a polytope $X$, take the cone over $\{1\}\times X$ in its affine span.
It is pointed, and the ray through $(1,v)$ is simple whenever $v$ is a
simple vertex. Homogeneous quadratic forms on this cone correspond to
quadratics on $X$. The cone is nonnegative exactly when the corresponding
quadratic is nonnegative on $X$. The moment-hull assertion follows from
compact homogenization and duality as in \eqref{repr:duals}.
\end{proof}

Every vertex of the cube is simple. The corollary does not
give a classification of polytopes without a simple vertex.

\subsection{Proper faces of the four-variable hull}

The obstruction in Theorem~\ref{repr:threshold} cannot be exhibited by
finding a proper face of $\mathcal Q_4$ without a finite lift.

\begin{proposition}\label{repr:proper-faces}
Every proper face of $\mathcal Q_4$ is a spectrahedral shadow.
\end{proposition}
\begin{proof}
An exposed proper face is the moment hull of the cube zero set of a nonzero
quadratic $p\geq0$. On the relative interior of a cube face containing a zero,
the restriction of $p$ has zero gradient and a positive semidefinite Hessian.
Its zero set on that face is therefore the intersection of the face with an
affine subspace. These intersections form a finite union of polytopes of
dimension at most three: cube boundary faces already have that dimension,
and a nonzero quadratic cannot vanish on a four-dimensional open set.
Their moment hulls have finite lifts by triangulation. The convex hull of
finitely many compact sets with finite lifts also has a finite lift, using
their closed homogenized cones and weights summing to one.
Thus every exposed proper face has a finite lift.

For a general proper face, take a supporting hyperplane at a point in its
relative interior. The resulting exposed face contains it. Repeating this
operation inside the containing face reaches the given face after finitely
many dimension drops. Affine sections preserve finite lifts.
Appendix~\ref{repr:face-proof} supplies the zero-set and face-chain details.
\end{proof}
